\par\bigskip\noindent\textbf{K101 supporting proof: The complete source, point-line, normal and selected-collar maps}\par\nobreak
Initial exact coordinate stage. The full line/root/cap attachment and
return remain to be computed; no return or compact boundary is assigned
here. This note retains the actual two parameter maps before any use of
a global loop relation.

\par

\medskip\noindent\textbf{EI.1. Source rescaling and selected normalization}

\par\nobreak

Put \(\epsilon=w^{-1}\), \(d_0=1-3\epsilon^2\).
In the marked source of EG.2 use
\[
 X=\epsilon^{-2}x,\qquad Y=\epsilon^{-3}y,\qquad
 T=\epsilon^{-2}t.
\]
Substitution gives the entire family
\[
 y^2=x^3+(-2+3\epsilon^2)x^2+d_0tx+\epsilon^2t^2,
 \qquad G=(0,\epsilon t).
 \tag{EI.1}
\]
The original E6 variables and parameter are
\[
 x_{\rm orig}=\frac{\epsilon^2}{d_0^2}x,\quad
 y_{\rm orig}=\frac{\epsilon^3}{d_0^3}y,\quad
 t_{\rm orig}=\frac{\epsilon^4}{d_0^3}t,\quad
 r=\frac{\epsilon^2(3\epsilon^2-2)}{d_0^2}.
 \tag{EI.2}
\]
Thus the original parameter has exact local order two, with retained
unit \((3\epsilon^2-2)/d_0^2\). The selected height is
\[
 t_c=1-\epsilon^2,\quad q_\epsilon=t-(1-\epsilon^2),\qquad
 q_{\rm orig}=\frac{\epsilon^4}{d_0^3}q_\epsilon.
 \tag{EI.3}
\]
Its whole normalization is
\[
 x=s^2-\epsilon^2,\quad y=s(s^2-1),\qquad
 Z=\frac{s-1}{s+1},\quad
 P_G=-\frac{1+\epsilon}{1-\epsilon}.
 \tag{EI.4}
\]
The ordered branches are \(s=-1,+1\), while \(G\) is \(s=-\epsilon\).
The selected Tate height coefficient in this chart is
\[
 \kappa_\epsilon(0,\epsilon)
       =-\frac{(1-\epsilon^2)^2}{64}.
 \tag{EI.5}
\]
This follows by multiplying the full earlier coefficient
\(- (w^2-1)^2/(64w^6)\) by the exact height scale
\(\epsilon^{-2}\); its phase is retained.

The original-source identity under \(\epsilon\mapsto-\epsilon\)
has the literal new-coordinate arrow
\[
 (x,y,t,\epsilon)\longmapsto(x,-y,t,-\epsilon).
 \tag{EI.6}
\]
It sends \(s\mapsto-s\), hence \(Z\mapsto Z^{-1}\) and
\(P_G\mapsto P_G^{-1}\). The point-line and finite-root coefficients
of this arrow still require their full comparison; EI.6 alone
does not identify a completed endpoint.

\par

\medskip\noindent\textbf{EI.2. Actual Miller scaling to be retained}

\par\nobreak

Write \(\alpha=\epsilon/d_0\). Then the original variables in EI.2
are \(x_{\rm orig}=\alpha^2x\),
\(y_{\rm orig}=\alpha^3y\),
\(t_{\rm orig}=\epsilon\alpha^3t\).
The tangent slope at \(G\) is \(d_0/(2\epsilon)\); the actual
second point has
\[
 x(2G)=\frac{(1-\epsilon^2)(1+3\epsilon^2)}
                  {4\epsilon^2}.
 \tag{EI.7}
\]
Consequently the exact original tangent-over-vertical factor is
\[
 f_{2,\rm orig}
   =\alpha\,f_{2,\epsilon},\qquad
 f_{2,\epsilon}
   =\frac{y-\epsilon t-d_0x/(2\epsilon)}
          {x-(1-\epsilon^2)(1+3\epsilon^2)/(4\epsilon^2)}.
 \tag{EI.8}
\]
Its entire numerator and denominator must be retained through the
section chart at infinity. On ordinary finite \(x,y,t\) charts,
\(f_{2,\epsilon}/\epsilon\) extends with value \(2x\).
No generic-Picard substitution is made for the actual \(M_2\) or
\(M_6\) Cartier map. The full six-factor suspension, the selected
height factor EI.3, both root progressions and their finite
characters are the next obligations.

\par

\medskip\noindent\textbf{EI.3. The full normalized Miller functions and central source map}

\par\nobreak

Set \(A_0=3\epsilon^2-2\), \(C_0=(1-\epsilon^2)(1+3\epsilon^2)\)
and \(S_0=C_0/4\). On every finite source-base chart the exact
scaled tangent function is
\[
 \widehat f_2=\frac{f_{2,\epsilon}}{\epsilon}
 =\frac{\epsilon y-\epsilon^2t-d_0x/2}
              {\epsilon^2x-S_0}.
 \tag{EI.9}
\]
Its central value is \(2x\), including the actual factor two.
The whole third-point Miller function can be written without a
negative parameter power as
\[
 f_{3,\epsilon}
 =\frac{S_0x(\epsilon y-\epsilon^2t-d_0x/2)}
 {S_0(\epsilon y+\epsilon^2t-d_0x/2)-2\epsilon^4tx}.
 \tag{EI.10}
\]
This is the original tangent-and-chord function under EI.2;
on the generic central source it is exactly \(x\).

Here are all factors in the sixth-point calculation. Define
\[
 \begin{aligned}
 X_3&=\frac{2d_0t}{S_0}+\frac{4\epsilon^4t^2}{S_0^2},\\
 Y_3&=t\left(\frac{d_0^2}{S_0}-1+
          \frac{6d_0\epsilon^4t}{S_0^2}
          +\frac{8\epsilon^8t^2}{S_0^3}\right),\\
 L_3&=\frac{3\epsilon^4X_3^2+2A_0\epsilon^2X_3+d_0t}
                     {2Y_3},\\
 X_6&=L_3^2-A_0\epsilon^2-2\epsilon^4X_3.
 \end{aligned}
 \tag{EI.11}
\]
Then the actual coordinates and tangent are
\(x(3G)=\epsilon^2X_3\), \(y(3G)=\epsilon Y_3\),
\(\ell_{3G}=\epsilon^{-1}L_3\), \(x(6G)=\epsilon^{-2}X_6\).
The complete tangent-over-vertical factor is therefore
\[
 \frac{f_{2,3G,\epsilon}}{\epsilon}
 =\frac{\epsilon y-\epsilon^2Y_3-L_3x+
                         \epsilon^2L_3X_3}
             {\epsilon^2x-X_6},\qquad
 \widehat f_6:=\frac{f_{6,\epsilon}}{\epsilon}
       =f_{3,\epsilon}^2\frac{f_{2,3G,\epsilon}}{\epsilon}.
 \tag{EI.12}
\]
At the generic central point \(X_3=8t\), \(Y_3=3t\),
\(L_3=1/6\), \(X_6=1/36\), hence
\(\widehat f_6=6x^3\). These are identities of the actual
meromorphic functions and their Cartier divisors. The factors
are retained at their zero and pole loci; no unit extension
through those horizontal divisors is asserted.

The exact source translation by \(-G\) is
\[
 k=\frac{y+\epsilon t}{x},\qquad
 x'=k^2-A_0-x,\qquad y'=\epsilon t-kx'.
 \tag{EI.13}
\]
At \(\epsilon=0\) it becomes
\[
 {\cal T}_0(x,y,t)=\left(\frac t x,-\frac{ty}{x^2},t\right),
 \qquad {\cal T}_0^2=1.
 \tag{EI.14}
\]
The finite I2 resolution is also carried literally. In the
chart \(x=t a,\ y=t b\) its equation is
\[
 b^2=t a^3+A_0a^2+d_0a+\epsilon^2.
 \tag{EI.15}
\]
The exceptional conic is smooth on the whole parameter disc
because \(d_0^2-4A_0\epsilon^2=C_0\) is a unit. In the
other source chart \(t=xu,\ y=xv_0\), use
\[
 \begin{aligned}
 x&=v_0^2-A_0-d_0u-\epsilon^2u^2,\qquad
 D_-=d_0+2\epsilon v_0+2\epsilon^2u,\\
 x'&=uD_-,\qquad u'=\frac{x}{D_-},\qquad
 v'_0=\epsilon u'-v_0-\epsilon u.
 \end{aligned}
 \tag{EI.16}
\]
The denominator is a unit near each central intersection.
Near the marked \(G\) point in EI.15, the reciprocal chord
\[
 k^{-1}=\frac{b-\epsilon}{t a^2+A_0a+d_0}
 \tag{EI.17}
\]
is an actual coordinate in the target \(O\) chart. The denominator
is nonzero. The inverse translation replaces \(\epsilon\) by
\(-\epsilon\) in these chord formulas. These arrows extend both
translations through the finite source resolution and retain
the exact \(t,\epsilon\).

The central smooth point \(G=(0,0)\) has order two. For both
odd \(b=1,3\) the suspension source map is exactly EI.14.
Substituting the whole functions EI.9-\/-12 in the two actual
products of EG.7 gives
\[
 \left.\epsilon^{-6}B_{6,1,\epsilon}\right|_{\epsilon=0}
 =\left.\epsilon^{-6}B_{2,3,\epsilon}\right|_{\epsilon=0}
 =64t^3 .
 \tag{EI.18}
\]
For the first product this is
\((2x)^6\,6(t/x)^3/(6x^3)\). For the second it is the
six chronological factors with \(x,\,t/x,\,x,\,t/x\) and
their prescribed multiplicities. Thus the exact base rescaling
in the actual point frames is
\[
 B_{{\rm sel},{\rm orig}}
 =\epsilon^{-2}\frac{q_\epsilon}{t^3}B_{a,b,\epsilon}
 =\epsilon^4\,\widehat B,\qquad
 \widehat B=\epsilon^{-6}\frac{q_\epsilon}{t^3}B_{a,b,\epsilon},
 \qquad \widehat B|_{\epsilon=0}=64q_\epsilon.
 \tag{EI.19}
\]
The retained original scalar is therefore
\(\lambda_{\rm native}=\epsilon^4\lambda\). Its actual
\(m\)-th root is \(\epsilon^{4/m}R_\lambda\); its endpoint
is \(\zeta_m^4R_\lambda\). This root path is retained together
with the finite source/root endpoint, not treated as an
integer number of closed root-parameter loops.

\par

\medskip\noindent\textbf{EI.4. Both physical line frames, half endpoint and full Cox collars}

\par\nobreak

Let \(\sigma_a\) be the actual canonical meromorphic point section.
The source map carries it literally. The full normalized
invariant differential and its transported vanishing period
scale by the same factor, so \(Q,P,Z\) agree on the source
comparison annulus. Thus
\[
 V_{\rm pt}=f_{P^a}v_M,\qquad
 V_{\rm tensor}=f_P^a f_a^{-1}v_M
 \tag{EI.20}
\]
have the directed coefficient arrows
\[
 \begin{aligned}
 V_{{\rm pt},{\rm orig}}&=V_{{\rm pt},\epsilon},\\
 f_{a,{\rm orig}}&=\alpha^{a-1}f_{a,\epsilon},\qquad
 V_{{\rm tensor},{\rm orig}}
       =\alpha^{1-a}V_{{\rm tensor},\epsilon},\\
 V_{{\rm tensor},w}&=\epsilon^{a-1}V_{{\rm tensor},\epsilon}.
 \end{aligned}
 \tag{EI.21}
\]
The two complete Cox-character maps are consequently
\[
 \begin{aligned}
 q_{\rm orig}^h Z^u V_{{\rm tensor},{\rm orig}}^j
 &=\epsilon^{\,4h-(a-1)j}d_0^{-3h+(a-1)j}
       q_\epsilon^hZ^uV_{{\rm tensor},\epsilon}^j,\\
 q_w^h Z^u V_{{\rm tensor},w}^j
 &=\epsilon^{-2h+(a-1)j}
       q_\epsilon^hZ^uV_{{\rm tensor},\epsilon}^j.
 \end{aligned}
 \tag{EI.22}
\]
Every parameter power in these arrows is retained.

For the native fixed-edge chart the actual characters are
\(q=abc,\ Z=b,\ V=c\). Put
\(\sigma=\epsilon^4d_0^{-3}\), \(\chi=\alpha^{1-a}\).
The new-to-original full germ is
\[
 (b,a,c)\longmapsto
       (b,\sigma\chi^{-1}a,\chi c).
 \tag{EI.23}
\]
The normal physical weights are \((-1,+1)\), so its actual
positive-weight coefficient arrow is
\[
 (\overline a,c)\longmapsto
      (\overline\sigma\,\overline\chi^{-1}\overline a,\chi c).
 \tag{EI.24}
\]
In the \(\epsilon\) loop these two phases have windings
\(-(a+3)\) and \(1-a\), respectively. The determinant
winding is \(-2a-2\). These are the actual frame coefficients,
not a characteristic-pair calculation. The entire diagonal
germ EI.23 and its parameter derivative are retained on the
annulus; its parameter-dependent source-normal terms have not
been discarded.

The half endpoint EI.6 preserves the canonical point section.
Its action on the complete Miller function is
\(f_{a,\epsilon}(x,-y,t,-\epsilon)=-f_{a,\epsilon}(x,y,t,\epsilon)\)
for \(a=2,6\). This follows either from the explicit chord
formulas or from EI.21 and the original-source identity.
The full theta inversion therefore gives
\[
 \begin{array}{c|cc}
 &Z&V\\ \hline
 \text{point frame}&Z^{-1}&P^aV\\
 \text{tensor frame}&Z^{-1}&-P^aV .
 \end{array}
 \tag{EI.25}
\]
After the literal target source-period \(h\), with line
coefficient \(P^{-a}\), these become \((Q/Z,V)\) and
\((Q/Z,-V)\), respectively. The height-one period and the
extra physical scalar \(-1\) in the tensor frame are retained.
Two such half arrows compose to the identity on these source
and line charts. Their full finite-cap attachment is being
constructed separately.

On the selected normalization one can see the exact point/tensor
pole without an asymptotic substitute. Put
\[
 s_a=\frac{1+P^a}{1-P^a},\qquad
 f_{a,\epsilon}=\frac{(s+\epsilon)^a}{s-s_a}.
 \tag{EI.26}
\]
Then
\[
 \frac{f_P^a}{f_{a,\epsilon}f_{P^a}}
 =\begin{cases}
 (2\epsilon)^{-1},&a=2,\\
 [\epsilon(6+20\epsilon^2+6\epsilon^4)]^{-1},&a=6 .
 \end{cases}
 \tag{EI.27}
\]
This is precisely the scalar parity in EI.25. The full smooth-fibre
ratio has the same actual Cartier and theta definition EI.20;
its entire source-base unit remains in the annular square.

\par

\medskip\noindent\textbf{EI.5. Finite central specialization and the remaining infinity divisor}

\par\nobreak

On every compact finite \(t\)-chart the section \(2G\) lies in the
identity formal chart at \(O\), with the exact coordinate
\[
 \tau(2G)=
 \frac{2\epsilon S_0}{d_0S_0+2\epsilon^4t}.
 \tag{EI.28}
\]
The denominator is a unit. The formal multiplication map
therefore gives the actual section \(aG=[a/2](2G)\) for
\(a=2,6\), with \(\tau(aG)=a\epsilon+O(\epsilon^3)\).
In the \(O\) Cartier chart its canonical section is the
coefficient \((\tau-\tau(aG))/\tau\). Together with the
canonical section away from these two divisors, this is the
whole point-line atlas; at \(\epsilon=0\) its transition is
one. Thus both actual point lines specialize to the trivial
line on the finite central source, including its I2 resolution.

Their point-to-tensor maps still contain their actual vertical
parameter divisors: EI.9 and EI.12 give order one along the
generic native divisor \(\epsilon=0\). In the original frame,
EI.21 gives order \(a\). The limits are \(2x\) and \(6x^3\),
whose zeros and poles remain the actual central Cartier
divisors of the tensor powers of the two-torsion point line.
No parameter power is removed from the annular map.

At the source-base infinity write \(q_\infty=t^{-1}\),
\(x_0=q_\infty^2x,\ y_0=q_\infty^3y\). The central equation is
\[
 y_0^2=x_0^3-2q_\infty^2x_0^2+q_\infty^3x_0.
 \tag{EI.29}
\]
Its actual invariants are
\[
 c_4=q_\infty^3(-48+64q_\infty),\quad
 c_6=q_\infty^5(-576+512q_\infty),\quad
 \Delta=64q_\infty^9(q_\infty-1).
 \tag{EI.30}
\]
The equation is minimal, since the \(q_\infty^3x_0\) coefficient
prevents division of \(x_0,y_0\) by \(q_\infty^2,q_\infty^3\).
These are the III* minimal orders. The marked section is
\((x_0,y_0)=(0,0)\). In the \(O\) chart its moving even multiples
have \(\tau_0(aG)=\tau(aG)/q_\infty\); this retains the pole
in the two variables \((\epsilon,q_\infty)\).
The full exceptional-component specialization and the finite
root/Cartier graph at this corner are not supplied by the
finite-chart trivialization above. They remain in the separate
root-atlas construction. No completed return is assigned here.

\par

\medskip\noindent\textbf{EI.6. The base graph and the complementary smooth source chart}

\par\nobreak

All windings in this note refer to the positive local epsilon loop.
The original source-base coordinate and the native infinity coordinate
obey the exact equation
\[
 t_{\rm orig}q_{\infty,\epsilon}=\epsilon^4d_0^{-3}.
 \tag{EI.31}
\]
Choose the analytic fourth root of the unit \(d_0^{-3}\) with value one.
Then \(\widetilde\epsilon=\epsilon d_0^{-3/4}\) identifies this base
graph with \(t_{\rm orig}q_{\infty,\epsilon}=\widetilde\epsilon^4\).
This is a directed analytic unit change, with its displayed inverse.
The original finite cap is the chart with
\(q_{\infty,\epsilon}=\widetilde\epsilon^4q_{\infty,\rm orig}\)
and bounded small \(q_{\infty,\rm orig}\). No enlargement of that cap
domain is part of this statement.

Use the complementary graph cover of EIR.7. For the two rows put
\(b_*=l\) when \(m=3l\), and \(b_*=l+2\) when \(m=3l+2\).
Its exact coordinates are
\[
 z=\xi^4,\quad \epsilon=\xi^m\sigma,\quad
 X_{\rm root}=\xi^{2b_*}\bar X,\quad
 Y_{\rm root}=\xi^{3b_*}\bar Y.
 \tag{EI.32}
\]
Both equations become
\[
 \bar Y^2=\bar X^3+a_2\bar X^2+d\bar X+\sigma^2,
 \quad a_2=A_0(\epsilon)\xi^{2m},\quad d=d_0(\epsilon),
 \qquad \bar G=(0,\sigma).
 \tag{EI.33}
\]
The original source map is
\[
 (x_o,y_o,t_o,r)=
 \left(\frac{\sigma^2}{d^2}\bar X,
       \frac{\sigma^3}{d^3}\bar Y,
       \frac{\sigma^4}{d^3},
       \frac{\xi^{2m}\sigma^2A_0(\epsilon)}{d^2}\right).
 \tag{EI.34}
\]
Substitution into the original cubic is literal. At \(\sigma=0\)
the curve is
\(\bar Y^2=\bar X^3-2\xi^{2m}\bar X^2+\bar X\), whose
discriminant is \(64(\xi^{4m}-1)\); it is a smooth compact elliptic
curve for the chosen small xi disc. The point \(\bar G=(0,0)\)
is its nonzero two-torsion point. Moreover \(t_o=0\) on this
divisor. This is the ordinary original III source chart; the original
cap domain, with \(q_{\infty,\rm orig}\) small, is at large sigma
in this cover and remains on its original K96 chart.

\par

\medskip\noindent\textbf{EI.7. The complete point-line, tensor and finite-deck arrows}

\par\nobreak

Put \(S=d^2/4-a_2\sigma^2=C_0(\epsilon)/4\). The actual second
section has
\[
 \bar X(2G)=S/\sigma^2,\qquad
 \tau(2G)=\frac{2\sigma S}{dS+2\sigma^4},\qquad
 \tau=-\bar X/\bar Y.
 \tag{EI.35}
\]
The denominator is a unit on the full small joint bidisc. The actual
formal multiplication map carries this section to
\(aG=[a/2](2G)\), for \(a=2,6\). In the O chart the canonical
point section has coefficient
\((\tau-\tau(aG))/\tau\), and away from those two divisors it
is the canonical meromorphic section. These are the full Cartier
transition maps. Their central transition is one, so they give a
literal central trivialization of \(M_a=\mathcal O(aG-O)\).

The entire tangent factor, with its horizontal divisor retained, is
\[
 \widehat f_{2}=\frac{\bar f_2}{\sigma}
  =\frac{\sigma\bar Y-\sigma^2-d\bar X/2}
             {\sigma^2\bar X-S}.
 \tag{EI.36}
\]
The full third and sixth factors are the formulas EI.10-\/-12 with
\((\epsilon,t,A_0,S_0)\) replaced by
\((\sigma,1,a_2,S)\). Thus, on the generic central curve,
\[
 \widehat f_2=2\bar X,\qquad
 \bar f_3=\bar X,\qquad
 \widehat f_6:=\bar f_6/\sigma=6\bar X^3.
 \tag{EI.37}
\]
The formulas retain the entire Cartier factors at O and G, rather than
claiming that the displayed powers of X are units there.

Let \(v_M\) be the coefficient in the canonical point section and
write \(V_{\rm tens,bar}=\bar f_a^{-1}v_M\). The vertically primitive
tensor coefficient and the actual original tensor coefficient are
\[
 W=\widehat f_a^{-1}v_M=\sigma V_{\rm tens,bar},\qquad
 f_{a,o}=\sigma^a d^{1-a}\widehat f_a,\qquad
 V_{{\rm tens},o}=\sigma^{-a}d^{a-1}W.
 \tag{EI.38}
\]
This is an exact morphism with its full vertical divisor. The
Cartier isomorphism using \(\widehat f_a\) carries the central point
trivialization to the tensor power of the two-torsion line, whose
canonical divisor is \(aG-aO\). Its displayed central functions are
\(2\bar X\) and \(6\bar X^3\). Equation EI.38, including its
sigma power, is used on every overlap with the original atlas.

The actual graph deck and the lift of the native root deck are
\[
 \begin{aligned}
 g_4:(\xi,\sigma,\bar X,\bar Y)&\longmapsto
  (i\xi,i^{-m}\sigma,i^{2m}\bar X,i^{3m}\bar Y),\\
 \widetilde D:(\xi,\sigma,\bar X,\bar Y)&\longmapsto
  (\zeta_{4m}\xi,-i\sigma,-\bar X,-i\bar Y).
 \end{aligned}
 \tag{EI.39}
\]
Here \(\widetilde D^m=g_4\), including all four coordinates.
Both arrows fix the original source functions EI.34. They carry
the actual point section to itself, since tau and tau(aG) transform
by the same scalar. Their full line coefficients are
\[
 \begin{array}{c|ccc}
 &v_M&\widehat f_a&W\\ \hline
 g_4&v_M&i^{ma}\widehat f_a&i^{-ma}W\\
 \widetilde D&v_M&i^a\widehat f_a&i^{-a}W .
 \end{array}
 \tag{EI.40}
\]
For either admitted even a, the first tensor character is
\((-1)^m\), and the second is \(-1\). Multiplication by the
literal \(\sigma^{-a}d^{a-1}\) in EI.38 makes the original
tensor coefficient invariant. Thus the central point character one
and the primitive tensor character in EI.40 are related by the
written vertical Cartier map; neither has been silently replaced.

These are also the full local normal maps. In the O chart tau
transforms by \(i^{-m}\) under \(g_4\), and by \(-i\) under
\(\widetilde D\); sigma has the same respective multiplier.
The corresponding local source-volume character on
\((\tau,\xi,\sigma)\) is \(i^{1-2m}\) or
\(-\zeta_{4m}\). The point-line coefficient is unchanged, so
the same determinant character holds after adjoining that complex
line. In the G chart, Y is a local source coordinate and gives
the identical characters. These complex-linear derivatives carry
the canonical complex spin-c structure; no choice of an unrecorded
square root or finite stabilizer lift is made.

The source translation by \(-G\) is the entire smooth-elliptic
translation in EI.33, with the chord and O-coordinate formulas
EI.13, EI.17, and their inverses. At sigma zero it is
\((\bar X,\bar Y)\mapsto(1/\bar X,-\bar Y/\bar X^2)\).
The full six-factor coefficient has vertical order six:
\[
 \bar B_{a,b}=\sigma^6\mathcal B_{a,b},\quad
 \mathcal B_{a,b}|_{\sigma=0}=64,
 \qquad B_{{\rm sel},o}=H\sigma^4\mathcal B_{a,b},
 \quad H=1-\xi^{4m}(1-\epsilon^2).
 \tag{EI.41}
\]
For completeness, the extension in this formula is an extension of
the actual line map. The source translation and the two point-line
Cartier atlases extend on the smooth proper elliptic family EI.33.
Consequently their Hom line extends. The displayed rational
six-factor section is a holomorphic isomorphism away from sigma
zero, so its divisor in this Hom line is supported on that single
irreducible central divisor. Its generic order is exactly six by
EI.36-\/-37 and the chronological factors; dividing by the displayed
sigma power gives divisor zero. In local Cartier frames both that
section and its inverse therefore extend holomorphically, and their
product is one. This proves the unit map \(\mathcal B_{a,b}\) on
the whole compact central curve, including O and G. The central
value 64 on a dense chart is consequently its value everywhere.
Here is the horizontal cancellation in those actual Cartier frames.
On this smooth joint family the complete divisors are
\[
 \begin{aligned}
 \operatorname{div}(\bar f_2)&=2[G]-[2G]-[O]+[\sigma=0],\\
 \operatorname{div}(\bar f_6)&=6[G]-[6G]-5[O]+[\sigma=0].
 \end{aligned}
 \tag{EI.41a}
\]
Pulling these divisors by the actual translations in the chronological
products gives
\[
 \operatorname{div}(\bar B_{a,b})
 =[aG]+[bG]-[(a+b)G]-[O]+6[\sigma=0].
 \tag{EI.41b}
\]
For \((a,b)=(6,1)\) this is the divisor of
\(\bar f_2^6\bar f_6(Tp)/\bar f_6(p)\). For \((2,3)\)
the four factors with exponents \(1,1,2,2\) at
\(p,T^3p,Tp,T^2p\) give exactly the same displayed formula.
Let \(r_p=(\tau-\tau(aG))/\tau\) in the O chart and
\(r_p=1\) in its canonical-section complement; the nonzero
point-line frame is \(e_p=\sigma_a/r_p\). In these frames the
whole coefficient, including both horizontal endpoint charts, is
\[
 \sigma^{-6}\bar B_{a,b}(p)
          \frac{r_{T_bp}}{r_p},\qquad T_b=t_{-bG}.
 \tag{EI.41c}
\]
Equations EI.41a-\/-b show its divisor is zero. This verifies the
extension at O, aG, bG and (a+b)G with the written local functions,
including their intersections at sigma zero. The inverse coefficient
has the opposite zero divisor, giving the inverse extending map.
The remaining factor \(H\sigma^4\), the old cap boundary domain,
and every intermediate graph component are retained. A compact
period quotient or a source return is not assigned by EI.41.

\par

\medskip\noindent\textbf{EI.8. The actual selected normalization graph in the original III resolution}

\par\nobreak

Put \(\nu=\epsilon^2/d_0\), choose the analytic square root of
\(d_0\) with value one, and set \(\rho=\epsilon/\sqrt{d_0}\).
Thus \(\rho^2=\nu\), while
\[
 r=-2\nu-3\nu^2,\qquad t_{c,o}=\nu^2(1+2\nu).
 \tag{EI.42}
\]
The original source resolution is the blowup of
\((x_o,y_o,t_o)\). Its chart \(t_o=x_ou,\ y_o=x_ov\) has
\(x_o=v^2-r-u-u^2\); the other chart has
\(x_o=t_o a_s,\ y_o=t_o b_s\) and
\(b_s^2=t_o a_s^3+r a_s^2+a_s+1\).
Their displayed coordinate changes are the actual source blowup maps.
The exceptional conic is smooth, since \(1-4r\) is a unit.

Use the original selected-normalization coordinate
\(S=(\epsilon/d_0)s\). Its exact map is
\[
 x_o=S^2-\nu^2,\quad
 y_o=S\{S^2-\nu(1+3\nu)\},\quad
 t_o=\nu^2(1+2\nu).
 \tag{EI.43}
\]
The pullback of the source blowup ideal is exactly
\[
 (x_o,y_o,t_o)=(S^2,S\nu,\nu^2)=(S,\nu)^2.
 \tag{EI.44}
\]
Indeed the last generator first supplies \(\nu^2\), the first
then supplies \(S^2\), and
\(y_o-S(S^2)+3S\nu^2=-S\nu\). Conversely all three
displayed source functions lie in the right-hand ideal. Taking
the second Veronese subalgebra of the Rees algebra gives the
literal Proj isomorphism
\(\operatorname{Bl}_{(S,\nu)^2}\cong\operatorname{Bl}_{(S,\nu)}\).
This proves the graph identification with the full source blowup,
including its y coordinate.

In the chart containing the two moving node branches put
\[
 h=-\nu/S,\qquad k=-S,
 \qquad hk=\rho^2.
 \tag{EI.45}
\]
Write \(u=\nu+U\). The actual graph-to-source map is
\[
 L=\frac{(1+3hk)h-k}{1-h^2},\qquad
 U=hL,\quad v=L,
 \tag{EI.46}
\]
and the source height is identically
\[
 q_o=(\nu+U)v^2-U^2(1+3\nu+U).
 \tag{EI.47}
\]
Substitution of EI.46 makes this height zero. This keeps the
\(A_1\) graph \(hk=\rho^2\) itself. Its double branches satisfy
\(k=h/(1-3h^2)\), with the two actual lifts of rho retained;
they map to \(U=v=0\). On the central axes, the h axis maps
to the exceptional source component, while the k axis maps to
the other III component. These claims are the displayed
polynomial/rational maps, not a branch relabelling.

The other graph chart is \(S=\nu z\), with
\[
 a_s=\frac{d_0(z^2-1)}{1-\epsilon^2},\qquad
 b_s=\frac{z(\epsilon^2z^2-1)}{1-\epsilon^2}.
 \tag{EI.48}
\]
At epsilon zero it maps the entire exceptional conic by
\((a_s,b_s)=(z^2-1,-z)\), including the marked G point
\(z=-1\). On the overlap \(h=-1/z\) and \(k=-\nu z\),
so EI.46 and EI.48 are literally the same source map.

Two ordinary blowups resolve the graph: first use
\(S=\rho z_1\), and then blow up \((z_1,\rho)\).
The charts \(z_1=\rho z\) and \(\rho=z_1\zeta\)
give \(S=\rho^2z\) and
\[
 (h,k,\rho)=(-\zeta,-z_1^2\zeta,z_1\zeta).
 \tag{EI.49}
\]
The first exceptional curve has self-intersection minus two
after the second blowup, and EI.49 contracts it, \(\zeta=0\),
to the graph point \(h=k=\rho=0\). The second exceptional
curve has self-intersection minus one and maps to the h-axis
component. This retains the actual contracted curve in the
source projection; no smooth embedded fixed trace is inferred
from the resolution alone.

The point and tensor lines on this graph also have an exact
description. For the admitted integer \(a=2,6\), let
\[
 S_a=\frac{\epsilon}{d_0}\frac{1+P^a}{1-P^a},
 \quad P=-\frac{1+\epsilon}{1-\epsilon},\qquad
 f_{a,o}=\frac{(S+\nu)^a}{S-S_a}.
 \tag{EI.50}
\]
The complete selected Miller formula EI.26 gives this identity;
\(S_a\) is analytic with \(S_a(0)=-1/a\). Thus the actual
point section of \(M_a\), whose divisor is the section \(aG\)
minus O, is a unit near this graph centre. The denominator in
EI.50 stays a unit there. Its tensor morphism retains the
numerator \((S+\rho^2)^a\). On the two exceptional valuations
\((\operatorname{ord}S,\operatorname{ord}\rho)=(1,1),(2,1)\)
the orders are exactly \(a,2a\). In the second chart it is
\(\rho^{2a}(z+1)^a\), retaining the strict marked G divisor;
the first chart retains \(\rho^a(z_1+\rho)^a\).
These line maps do not assign a characteristic contribution
to either exceptional curve.

\par

\medskip\noindent\textbf{EI.9. The ordinary finite quotient and its H-unit line chain}

\par\nobreak

On the complementary chart, the fourth power of
\(\widetilde D\) acts only by
\(\xi\mapsto\zeta_m\xi\), fixing sigma, the source point,
and its genuine point-line coefficient. Its invariant coordinate
is \(\Xi=\xi^m\). The invariant coordinate algebra therefore
gives the exact ordinary quotient chart
\[
 \bar Y^2=\bar X^3+A_0(\epsilon)\Xi^2\bar X^2
                 +d_0(\epsilon)\bar X+\sigma^2,
 \qquad \epsilon=\Xi\sigma.
 \tag{EI.51}
\]
The remaining deck is
\[
 (\Xi,\sigma,\bar X,\bar Y,v_M)
 \longmapsto(i\Xi,-i\sigma,-\bar X,-i\bar Y,v_M).
 \tag{EI.52}
\]
The inverse image on the native source is
\(q_{\infty,\epsilon}=\Xi^4\),
\(x=\Xi^{-2}\bar X\), \(y=\Xi^{-3}\bar Y\).
The quotient, its ramification, and its remaining fourth-root
action are written explicitly; the finite cap on its original
domain retains its original m-root and characters.

Here is a local partial filling on a closed small Xi disc, where
\(H=1-\Xi^4(1-\epsilon^2)\) is a unit. The actual source
translation is \(T=t_{-bG}\), and the forward suspension is
\[
 (p,V)\longmapsto
 (Tp,\lambda H\sigma^4\mathcal B_{a,b}(p)V).
 \tag{EI.53}
\]
For each integer j use the actual M and inverse-M coefficient
charts
\[
 X_jY_j=\sigma,\qquad V=\sigma^jX_j.
 \tag{EI.54}
\]
On a change of point-line frame, X and Y receive the mutually
inverse full transition functions of M. Adjacent charts glue
on \(Y_j\ne0\) by
\(X_{j+1}=Y_j^{-1}\), \(Y_{j+1}=\sigma Y_j\).
This is the full line chain over the same source, not a scalar
substitution in one punctured chart.

Put \(C=\lambda H\mathcal B_{a,b}\), the actual Hom-line
isomorphism over T. Equation EI.53 is exactly
\[
 B:(p,j,X_j,Y_j)\longmapsto
       (Tp,j+4,C(p)X_j,C(p)^{-1}Y_j).
 \tag{EI.55}
\]
The remaining fourth-root deck acts by
\[
 j\longmapsto j,\quad
 X_j\longmapsto i^jX_j,\quad
 Y_j\longmapsto(-i)i^{-j}Y_j,
 \tag{EI.56}
\]
together with EI.52 on the source. These two arrows commute
literally: the Hom section in EI.53 is the original invariant
section, and changing j by four does not change either
character in EI.56. All graph and line overlaps are therefore
retained, including their finite stabilizers.

The cover chain is smooth: the source is a smooth family over
the two base coordinates, and the equation in EI.54 has
nonzero sigma derivative at its nodes. At sigma zero B shifts
the chain index by four. On a compact source patch its Hom
norm is bounded; for small nonzero sigma the actual line norm
is contracted by at most \(C_1|\lambda||\sigma|^4\).
Thus finite consecutive chain regions and their literal B
identifications give a proper local quotient, with compact
central fibre. The subsequent finite fourth-root quotient is
retained with its actual stabilizers; smoothness of that
coarse quotient is not asserted here. The original m-cap
attaches only along its original-domain overlap EIR.7. The
boundary at \(H=0\), which is the selected colliding-node
region, still needs its full conifold/fan attachment. No
global compact filling or return is inferred from this
H-unit partial chain.

\par

\medskip\noindent\textbf{EI.10. The compact O and T section graphs and their cyclic reset}

\par\nobreak

The map from EI.51 to the original resolved source has two section
graphs on sigma zero. They are required before using the original
point-line frames on that resolution. At O put
\(\eta=-\bar X/\bar Y\), \(U=\eta^2\bar X\). The unit U is
specified by the exact equation
\[
 1=U+a_2\eta^2+d\eta^4/U+\sigma^2\eta^6/U^2,
 \qquad U(0)=1.
 \tag{EI.57}
\]
Blow up \((\eta,\sigma)\). In \(\eta=\sigma p\), the original
O coordinate is \(\tau_o=dp\). In \(\sigma=\eta v\) the whole
original source and its resolved coordinates are
\[
 x_o=v^2U/d^2,\quad y_o=-v^3U/d^3,\quad
 t_o=\eta^4v^4/d^3,\qquad
 u_o=\eta^4v^2/(dU),\quad v_o=-v/d.
 \tag{EI.58}
\]
All these expressions are regular. In the first chart the two
projective source coordinates are also regular, with \(\tau_o=dp\)
at p zero; the inverse source chart uses \(v=p^{-1}\).
Thus both halves of the section graph carry the full projective
source map, rather than only its affine part.

Let \(c_a=\tau(aG)/\sigma\), the analytic unit defined by the
actual formal multiplication in EI.35; \(c_a(0)=a\). The whole
point-section coefficient on this graph is
\[
 \frac{\eta-\tau(aG)}{\eta}
       =\frac{p-c_a}{p}=1-c_av.
 \tag{EI.59}
\]
The two total transforms have the same exceptional multiplicity,
which cancels in their difference. Thus the pullback of the actual
point-line atlas is precisely the strict point divisor aG minus O.
No vertical point-line correction has been inserted. The finite deck
fixes p and v, sends \((\Xi,\sigma)\) to
\((i\Xi,-i\sigma)\), and preserves \(c_a\) and EI.59 literally.

At T, the unit \(D_T=d+a_2\bar X+\bar X^2\) gives
\(\bar X D_T=\bar Y^2-\sigma^2\). Consequently the pullback
of the original source blowup ideal, after its common scalar factor,
is exactly
\[
 (\bar X,\sigma\bar Y,\sigma^2)
       =(\bar Y^2,\sigma\bar Y,\sigma^2)
       =(\bar Y,\sigma)^2.
 \tag{EI.60}
\]
The quadratic Veronese therefore gives the literal graph blowup
\((\bar Y,\sigma)\). In its first chart
\(\bar Y=\sigma z\), \(\bar X=\sigma^2 A\), the equation and
the original exceptional conic map are
\[
 z^2=\sigma^4A^3+a_2\sigma^2A^2+dA+1,
 \qquad (x_o/t_o,y_o/t_o)=(dA,z).
 \tag{EI.61}
\]
In the other chart \(\sigma=\bar Yv\), put
\(\bar X=\bar Y^2C\), where
\(1=\bar Y^4C^3+a_2\bar Y^2C^2+dC+v^2\).
Then \((u_o,v_o)=(v^2/(dC),v/(dC))\) wherever C is a unit;
the first chart covers its two marked zeros. The actual point line
has its canonical nonzero frame here because aG and O lie in the
other section chart. The finite deck fixes z or v and again acts
on the two transverse variables by \((i,-i)\).

The source translation permutes these two graph centres and hence
lifts to their blowups. This lift is also explicit. The chord
\(k=-1/\eta+\sigma\eta^2/U\) gives, on \(\eta=\sigma p\),
\[
 \frac{\bar X'}{\sigma^2}
  =\frac{dp^2-2p}{U}+\frac{2\sigma^4p^4}{U^2},\qquad
 \frac{\bar Y'}{\sigma}
  =1+\frac{\bar X'}{\sigma^2p}
       -\frac{\sigma^2p^2\bar X'}{U}.
 \tag{EI.62}
\]
At sigma zero these are \(A=p^2-2p\), \(z=p-1\).
The inverse source chart gives p equal to z minus one for the next
forward translation. Thus for the actual odd b, including every
intermediate primitive translation,
\[
 O:p\longmapsto T:z=p-b,\qquad
 T:z\longmapsto O:p=z-b,
 \qquad B^2:p\longmapsto p-2b.
 \tag{EI.63}
\]
The central coefficient of the Hom unit is 64 in the pulled-back
nonzero point frames. In the canonical meromorphic coefficients,
EI.41c and EI.59 give the actual two rational functions
\[
 64\frac{p-a}{p}\quad\hbox{on O},\qquad
 64\frac{z-b}{z-b-a}\quad\hbox{on T}.
 \tag{EI.64}
\]
This retains all four horizontal point divisors. The original
coefficient still includes \(\lambda H\sigma^4\).

Finally, the other blowup chart has \(\sigma=\eta v\) at O
and \(\sigma=\bar Yv\) at T. Pulling back the line chain EI.54
therefore gives the associated node \(X_jY_j=\eta v\), or
\(X_jY_j=\bar Yv\). The pre-graph smoothness argument in EI.9
does not erase these nodes. Their actual source-reset arrow,
finite characters, line map and relative resolution remain part
of the ordinary-to-selected attachment.

\par

\medskip\noindent\textbf{EI.11. A six-chart selected collar with its exact T0 boundary}

\par\nobreak

Near any one of the four solutions \(\Xi_0^4=1\), the derivative
of \(H=1-\Xi^4(1-\Xi^2\sigma^2)\) with respect to Xi is nonzero.
Use the actual local inverse \(\Xi=\Xi(\sigma,H)\). The marked
selected Tate coordinate is \(q=Q=\kappa(\sigma,H)H\), with the
entire nonzero kappa. The physical point-theta coefficient gives
the exact periods
\[
 h:(Z,V)\mapsto(qZ,\alpha V),\qquad
 B:(Z,V)\mapsto(gZ,\sigma^4q C V),
 \tag{EI.65}
\]
where \(\alpha=P^a\), \(g=P^{-b}\), and
\[
 C=\lambda\beta/\kappa,\qquad
 \beta=\sigma^{-6}\bar B_{a,b}(p)
       \frac{f_{P^a}(P^{-b}Z)}{f_{P^a}(Z)}.
 \tag{EI.66}
\]
The direction in EI.66 is the actual coefficient rule
\(V=f_{P^a}v_M\). Its source divisor cancels by EI.41b and the
full theta ratios. It is therefore a nonzero function of the base,
with central value 64; it can be evaluated on any fixed regular
Z chart, retaining the whole function. In particular C is a unit
on the specified collar. No logarithmic constant has been discarded.

On the universal source and line chains put
\[
 x_i=Z/q^i,\quad y_i=q^{i+1}/Z,\qquad
 U_j=V/(\sigma^4q)^j,\quad
 V_j=(\sigma^4q)^{j+1}/V.
 \tag{EI.67}
\]
Their common local equation is
\[
 U_jV_j=\sigma^4x_i y_i,\qquad q=x_i y_i.
 \tag{EI.68}
\]
Resolve its ordered factors
\((f_1,\ldots,f_6)=(\sigma,\sigma,\sigma,\sigma,y_i,x_i)\)
by the six actual charts
\[
 f_r=s_rt_r,\quad
 U_j=s_r\prod_{k<r}f_k,\quad
 V_j=t_r\prod_{k>r}f_k.
 \tag{EI.69}
\]
The remaining independent base variables give four complex
coordinates in each chart. For a sigma factor they are
\((s_r,t_r,x_i,y_i)\); for y or x they are the analogous
\((s_r,t_r,\sigma,x_i)\) or
\((s_r,t_r,\sigma,y_i)\). Thus every chart is smooth, with
physical weights \((+1,-1)\) on \((s_r,t_r)\).

In the lattice \((\operatorname{ord}Z,\operatorname{ord}V,
\operatorname{ord}\sigma,\operatorname{ord}H)\), put
\(E_j=(0,j,1,0)\), \(I_n=(n_1,n_2,0,1)\).
The six cones are
\[
 \begin{gathered}
 (E_{r-1},E_r,I_0,I_{e_1})\quad(1\le r\le4),\\
 (E_4,I_0,I_{e_1},I_{e_2}),\qquad
 (E_4,I_{e_1},I_{e_2},I_{e_1+e_2}).
\end{gathered}
 \tag{EI.70}
\]
Their primitive determinants have absolute value one. At nonzero
sigma the final two triangles are exactly
\(\{0,e_1,e_2\}\) and \(\{e_1,e_2,e_1+e_2\}\).
This is the actual negative-square T0 marking.

The full overlaps are retained. On \(y_i\ne0\), the next source
chart has \(x'=1/y_i\), \(y'=x_i y_i^2\). The sigma-factor
charts retain s and t; the old last x-factor chart changes to
the new penultimate y-factor chart by
\(s'=y_i s,\ t'=y_i t\), which preserves U and V literally.
The line-chain overlap \(U_{j+1}=V_j^{-1}\),
\(V_{j+1}=\sigma^4q V_j\) lies in the corresponding smooth
unit locus. These are genuine chart maps on their specified
overlaps.

For any diagonal unit arrow \(f_k'=a_kf_k\),
\(U'=cU\), \(V'=dV\), with \(cd=\prod a_k\), its full lift is
\[
 s_r'=c\Bigl(\prod_{k<r}a_k\Bigr)^{-1}s_r,\qquad
 t_r'=d\Bigl(\prod_{k>r}a_k\Bigr)^{-1}t_r.
 \tag{EI.71}
\]
It preserves EI.69 with all normal derivatives and units. For h,
the labels change by \((i,j)\mapsto(i+1,j)\), with
\((s,t)\mapsto(\alpha s,\alpha^{-1}t)\). For B the labels
change by \((i,j)\mapsto(i,j+1)\); the sigma charts have factors
\((C,C^{-1})\), the y chart \((C,C^{-1}g^{-1})\), and the x
chart \((Cg,C^{-1})\). These commute by their original period
maps. The fourth-root deck permutes the four H collars freely,
fixes Z, q and V, and sends sigma to minus i times sigma. Its
lift on the first four charts is
\((s_r,t_r)\mapsto(i^{r-1}s_r,i^{4-r}t_r)\), and it fixes
the s,t coordinates in the final two charts. All finite arrows
therefore retain the actual T0 boundary.

This gives the complete local selected collar and its period
actions. To combine it with EI.10, a section-corner chart has
\(\sigma=\eta v\); the corresponding line equation is
\(UV=\eta^4v^4q\). Its ordered factors
\((\eta,\eta,\eta,\eta,v,v,v,v,q)\) give nine smooth relative
charts by EI.69. On either unit overlap they reduce to the
four-sigma plus q charts. The actual source reset carries eta
and v to eta times a unit and v divided by the same unit near
their intersection (its central O-to-T factor is \(1-v\));
EI.71 retains that entire unit. The compact fundamental-region
and full old-cap matching for the union of these ordinary and
selected graph charts are still being audited. No global
characteristic value is assigned in this note.

The overlap with the H-unit ordinary line chain includes a full
integer monomial. In the r-th sigma chart of EI.69,
\[
 V_{\rm phys}=(\sigma^4q)^jU_j
       =q^j\sigma^{4j+r-1}s_r.
 \tag{EI.72}
\]
Consequently its ordinary chain index is \(k=4j+r-1\), with
\(X_k=q^js_r\), \(Y_k=q^{-j}t_r\). The B arrow sends k
to k plus four and has multiplier \(qC=\lambda H\beta\),
as required by EI.55 in the same point-theta frame. The entire
integer power \(q^j\) remains in the annular and finite-cap
comparison.

At the section-corner graph, put \(\eta=-\bar X/\bar Y\),
\(\sigma=\eta v\), \(U=\eta^2\bar X\). The full next source
chart from the chord is
\[
 \bar X'=\eta^2 A,\quad \bar Y'=\eta c,\qquad
 A=(d-2v)/U+2\eta^4v^2/U^2,\quad
 c=v+A-\eta^4vA/U.
 \tag{EI.73}
\]
Thus \((\eta,v)\mapsto(\eta c,v/c)\), with
\(c|_{\eta=0}=1-v\). For the actual odd b the chronological
composition gives its full unit \(c_b\), with central value
\(1-bv\). Near the retained corner this is nonzero. The nine
factors transform by \((c_b,c_b,c_b,c_b,
c_b^{-1},c_b^{-1},c_b^{-1},c_b^{-1},1)\). Equation EI.71
is the literal lifted map, including its source derivatives.

\par

\medskip\noindent\textbf{EI.12. Full ordinary, native and selected frame overlaps}

\par\nobreak

On \(\Xi\ne0\), the exact native source arrow is
\[
 (x_\epsilon,y_\epsilon,t_\epsilon)
   =(\Xi^{-2}\bar X,\Xi^{-3}\bar Y,\Xi^{-4}),
 \qquad \epsilon=\Xi\sigma.
 \tag{EI.74}
\]
It retains the full Miller factors
\(f_{a,\epsilon}=\Xi^{1-a}\bar f_a\) and
\(B_{{\rm raw},\epsilon}=\Xi^{-6}\bar B\).
Since \(q_\epsilon=\Xi^{-4}H\), the actual entire coefficient
satisfies
\[
 \widehat B_\epsilon
 =\epsilon^{-6}\frac{q_\epsilon}{t_\epsilon^3}
                     B_{{\rm raw},\epsilon}
 =\Xi^{-4}H\sigma^{-6}\bar B,\qquad
 \lambda\epsilon^4\widehat B_\epsilon
 =\lambda\sigma^4H\mathcal B.
 \tag{EI.75}
\]
The full point frames agree: \(\tau_\epsilon=\Xi\bar\tau\)
and \(\tau_\epsilon(aG)=\Xi\bar\tau(aG)\), so their
canonical Cartier ratios are identical. For native line charts
\(A_kB_k=\epsilon\), \(V=\epsilon^kA_k\), the whole ordinary
overlap is
\[
 X_k=\Xi^kA_k,\qquad
 Y_k=\Xi^{-k-1}B_k.
 \tag{EI.76}
\]
It preserves \(X_kY_k=\sigma\) and V literally. In particular
its integer powers reproduce the finite fourth-root characters
of EI.56. The finite native I2 resolution EI.15-\/-17 and its
point-line maps are transported by this exact arrow on the
specified overlap.

For the point-theta comparison in EI.65, use the nonzero
Cartier frame \(e_p=\sigma_a/r_p\) of EI.41c and put
\[
 \psi_p=\frac{f_{P^a}(Z(p))}{r_p}.
 \tag{EI.77}
\]
Both numerator and denominator retain the identical point and
O divisors; their quotient and its inverse are nonzero
holomorphic functions on the actual source overlap. This
includes the two section graphs, because both factors have
the same total transforms there. If v is the coefficient in
the nonzero frame, the complete arrow is \(V=\psi_p v\).
Thus the original-frame version of EI.72 is
\[
 X_k=\psi_p^{-1}q^js_r,\qquad
 Y_k=\psi_p q^{-j}t_r.
 \tag{EI.78}
\]
Every point-line unit, including its parameter dependence,
is retained by this formula. The inverse target source
factor is the one already displayed in EI.66 and EI.41c.

The source graph coordinates themselves have the same full
comparison. At the smooth O point the ratio
\((Z-1)/\bar\tau\) is an analytic unit; on the selected
normalization EI.4 its exact expression, with \(z=s^{-1}\), is
\[
 \frac{Z-1}{\bar\tau}
    =\frac{2\Xi(1-z)}{1-\epsilon^2z^2}.
 \tag{EI.79}
\]
Its value at O is \(2\Xi\). At the smooth central T point,
\((Z+1)/\bar Y\) is a unit with value \(-2\Xi^{-3}\).
The full smooth-source Tate maps retain those analytic unit
ratios, not just their values. Consequently their section
blowup ideals are respectively \((Z-1,\sigma)\) and
\((Z+1,\sigma)\), with the entire unit maps in both
directions. This identifies the source graph corners used
in EI.73 with the actual selected source charts. The
fourth-root phases are retained through the displayed Xi
factors. The original cap overlap still uses EIR.28 on its
original domain, with both signed norm choices, finite
characters, and the actual reduced-row B-power reindex.

For the full T-chart assertion, solve the cubic at \(\bar Y=0\)
for its unique root near \(\bar X=0\). This is the actual
two-torsion section \(T(\sigma,H)\). In the marked Tate chart
it is the nonzero two-torsion section whose limit is nonzero;
the two other sections have coordinates \(\pm\sqrt q\),
whereas this section is exactly \(Z=-1\). Thus \(Z+1\) is
divisible by \(\bar Y\) on the full smooth source chart,
and the quotient is a unit with the stated central value.
This proves \((Z+1,\sigma)=(\bar Y,\sigma)\) as actual
ideals. It does not replace the moving marked point
\(P_G(\sigma,H)\) by minus one away from sigma zero.

On the original finite-cap overlap the source angular lift is
\[
 s_o=\frac{3\log d}{2\pi i}-\frac{4\log\sigma}{2\pi i}.
 \tag{EI.80}
\]
The graph generator \(\sigma\mapsto-i\sigma\), with its
specified negative quarter-turn path, changes this lift by
plus one. The full K96 sector square therefore transports it
to \(A B^{j_\gamma}\), where \(j_\gamma\) is the prescribed
integral reference cochain. The full
\(A=R_\lambda^n\zeta_m^d\widetilde F^{\,n}D\) is retained.
Four iterates are \(A^4B^{\sum j_\gamma}\) on the sector
cover, with the literal norm \(A^m=B^n\). On the completed
cap quotient this gives the fourth-root graph action. The
ordinary-to-cap maps retain these B-decks; no zero value of
the reference integer is imposed.

